\section{System Architecture}
% Kürzen?
The architectural evolution examined in this paper is illustrated in \autoref{fig:workflow_vergleich} and contrasts two fundamental RAG designs. Version 1 (V1) corresponds to a naive RAG approach, while Version 2 (V2) represents an evolution toward an adaptive, partially modular system, as discussed in current research.

V1 follows a strictly linear pipeline consisting of contextualization, a semantic retrieval component, generation, and postprocessing. This design corresponds to the Naive RAG described in the literature, in which retrieved documents are integrated into the prompt unchanged. Although this approach is easy to implement, its performance depends heavily on retrieval quality and exhibits weaknesses with complex queries as well as when handling irrelevant contexts.

V2 addresses these limitations through an adaptive architecture with a central routing component. The Router \& Contextualizer classifies incoming queries and dynamically routes them to different processing paths. A retrieval-free path enables the efficient answering of simple or conversational queries without external knowledge retrieval, thereby reducing latency and potential noise. For knowledge-intensive queries, an extended pipeline is used that combines hybrid retrieval with subsequent reranking. This design decision reflects empirical findings indicating that retrieval quality is a key performance factor for RAG systems. Additionally, query decomposition is used to support multi-hop reasoning. The architecture thus moves toward Advanced RAG, but deliberately avoids fully modular or iterative systems. In particular, approaches such as GraphRAG were not integrated due to their high implementation effort and additional system complexity, even though they offer advantages for reasoning-intensive tasks. A key enhancement in V2 is the integration of a Socratic learning mode. This transforms the chatbot from a pure question-answering system into a didactically oriented assistance system that supports active knowledge construction through targeted questioning, thereby aligning with current findings in educational research. 

\begin{figure}[htbp]
\centering

\tikzset{
    box/.style={
        draw,
        rectangle,
        fill=white,
        text width=3.0cm,
        minimum height=1.15cm,
        align=center
    },
    line/.style={
        -Latex,
        thick
    }
}

% ===================== VERSION 1 =====================
\begin{minipage}{\textwidth}
\centering
\begin{tikzpicture}[
    font=\small,
    >=Latex
]

\node[box] (c1) at (1.8,0)  {Contextualizer};
\node[box] (r1) at (5.93,0)  {Semantic\\Retrieval};
\node[box] (g1) at (10.06,0)  {Generator};
\node[box] (p1) at (14.2,0) {Postprocessing};

\draw[line] (c1.east) -- (r1.west);
\draw[line] (r1.east) -- (g1.west);
\draw[line] (g1.east) -- (p1.west);

\end{tikzpicture}

\vspace{0.4em}
\textbf{(a) Version 1}
\end{minipage}

\vspace{1.2em}

% ===================== VERSION 2 =====================
\begin{minipage}{\textwidth}
\centering
\begin{tikzpicture}[
    font=\small,
    >=Latex
]

% Router
\node[box] (router) at (1.8,0) {Router \&\\Contextualizer};

% Split point
\coordinate (split) at (4.0,0);

% Upper workflow
\node[box] (gen1) at (14.2,1.8) {Generator};

% Middle workflow
\node[box] (retrieval) at (6.8,0)  {Hybrid\\Retrieval};
\node[box] (rerank)    at (10.5,0) {Reranker};
\node[box] (gen2)      at (14.2,0) {Generator \&\\Postprocessing};

% Lower workflow
\node[box] (socraticmode) at (14.2,-1.8) {Socratic Learning\\Mode};

% Trunk
\draw[thick] (router.east) -- (split);

% Branches from split
\draw[line] (split) |- (gen1.west);
\draw[line] (split) -- (retrieval.west);
\draw[line] (split) |- (socraticmode.west);

% Middle workflow
\draw[line] (retrieval.east) -- (rerank.west);
\draw[line] (rerank.east) -- (gen2.west);

% Labels at branching point
\node[anchor=west, align=left] at (4,2) {Path without Retrieval};
\node[anchor=west]             at (4,0.21) {Q\&A};
\node[anchor=west]             at (4,-1.57) {Socratic};

\end{tikzpicture}

\vspace{0.4em}
\textbf{(b) Version 2}
\end{minipage}

\caption{Conceptual workflows of the RAG systems evaluated in this work.}
\label{fig:workflow_vergleich}
\end{figure}

Beyond the core architecture, the knowledge base was significantly expanded through an improved ingestion pipeline, which nearly doubled the volume of retrievable tokens (from about 6 million to 11 million) and addressed shortcomings in the extraction of stock-related content. In addition, several measures were implemented to improve quality: First, response streaming reduces perceived latency through incremental output. Second, system prompts were revised in accordance with structured prompt engineering principles \autocite{Lee2025PromptEngineering}, leading to more consistent and controllable responses. Third, the citation system was adapted so that clear source attributions enhance the transparency and traceability of the generated content.